\documentclass[11pt]{article}
\usepackage[margin=1in]{geometry}
\usepackage{amsmath,amsfonts,amssymb,amsthm,enumerate,bbm}
\usepackage[]{graphicx}
\usepackage{color,subfigure}
\definecolor{scarlet}{rgb}{0.81,0,0}
\usepackage{multicol}
\usepackage{float}
\usepackage[all]{xypic}
\usepackage[colorlinks=true,citecolor=scarlet,linkcolor=scarlet]{hyperref}
\usepackage{colonequals}

\usepackage{fancyhdr, lastpage}
\pagestyle{fancy}
\fancyfoot[C]{{\thepage} of \pageref{LastPage}}



\DeclareMathOperator{\mSpec}{mSpec}
\DeclareMathOperator{\Spec}{Spec}
\DeclareMathOperator{\Ass}{Ass}
\DeclareMathOperator{\Supp}{Supp}
\DeclareMathOperator{\height}{height}
\DeclareMathOperator{\Hom}{Hom}
\DeclareMathOperator{\ann}{ann}
\DeclareMathOperator{\End}{End}
\DeclareMathOperator{\coker}{coker}
%\DeclareMathOperator{\ker}{ker}
\DeclareMathOperator{\rank}{rank}
\DeclareMathOperator{\im}{im}
\DeclareMathOperator{\M}{M}
\DeclareMathOperator{\Tor}{Tor}
\DeclareMathOperator{\id}{id}
\DeclareMathOperator{\ch}{char}
\DeclareMathOperator{\Aut}{Aut}

%\DeclareMathOperator{\dim}{dim}

\DeclareMathOperator{\lcm}{lcm}

\def\ra{\rightarrow}
\newcommand{\m}{\mathfrak{m}}
\newcommand{\C}{\mathbb{C}}
\newcommand{\Q}{\mathbb{Q}}
\newcommand{\Z}{\mathbb{Z}}
\newcommand{\ZZ}{\mathbb{Z}}
\newcommand{\R}{\mathbb{R}}
\newcommand{\N}{\mathbb{N}}
\newcommand{\ov}[1]{\overline{#1}}
\newcommand{\norm}{\trianglelefteq}

\def\ov#1{\overline{#1}}


\title{}
\date{\vspace{-0.5in}}

\makeatletter
\g@addto@macro\@floatboxreset\centering
\makeatother

\theoremstyle{definition}
\newtheorem{problem}{Problem}


\begin{document}

\thispagestyle{fancy}
\pagestyle{fancy}
\rhead{UNL $\mid$ Spring 2026}
\lhead{Introduction to Modern Algebra II}

\vspace{3em}

\begin{center}
	{\LARGE Problem Set 8 \\}
	Due Thursday, March 26
\end{center}

\

\noindent
{\bf Instructions:}
You are encouraged to work together on these problems, but each student should hand in their own final draft, written in a way that indicates their individual understanding of the solutions. Never submit something for grading that you do not completely understand. You cannot use any resources besides me, your classmates, and our course notes.


I will post the .tex code for these problems for you to use if you wish to type your homework. If you prefer not to type, please  {\em write neatly}. As a matter of good proof writing style, please use complete sentences and correct grammar. You may use any result stated or proven in class or in a homework problem, provided you reference it appropriately by either stating the result or stating its name (e.g. the definition of ring or Lagrange's Theorem). Please do not refer to theorems by their number in the course notes, as that can change.


\

\begin{problem} Let 
\[ A = \begin{bmatrix} -2 & 0 & 0 \\  -1 & -4 & 0 \\ 2 & 4 & 0 \end{bmatrix}  \qquad \text{and} \qquad B = \begin{bmatrix} -2 & 0 & 0 \\  -1 & -4 & -1 \\ 2 & 4 & 0 \end{bmatrix} \in \mathrm{Mat}_{3 \times 3}(\R).\]
\begin{enumerate}
\item[(a)] Find the rational canonical form for $A$ and $B$.
\item[(b)] Find the Jordan canonical form for $A$ and $B$, if they exist.
\item[(c)] Determine if $A$ and $B$ are diagonalizable.
\end{enumerate}
\end{problem}

\


\begin{problem}
Let $F$ be a field and $A,B\in \mathrm{Mat}_{n}(F)$.
\begin{enumerate}
\item[(a)] Show that $A$ is similar to its transpose $A^T$.
\item[(b)] Let $L\supseteq F$ be a field extension\footnote{That is, $F$ is a subfield of $L$.} of $F$. Show that if $A$ and $B$ are similar over $L$, then $A$ and $B$ are similar over $F$.
\end{enumerate}
\end{problem}



\

\begin{problem} List all possible rational canonical forms over $\Q$ and Jordan canonical forms over $\C$
for $8 \times 8$ matrices with determinant $81$ and minimal polynomial $(x-3)^2(x^2 + 1)$. Carefully justify your answer.
\end{problem}

\

\begin{problem} Let $V$ be a finite-dimensional vector space over a field $F$. Let $\phi:V\to V$ be a linear transformation. For an eigenvalue $\lambda$ of $\phi$, we say that
\begin{itemize}
\item the \textbf{arithmetic multiplicity} of $\lambda$ is the multiplicity of $\lambda$ as a root of $c_\phi(x)$; i.e., the largest $e$ such that $(x-\lambda)^e \ | \ c_\phi(x)$, and
\item the \textbf{geometric multiplicity} of $\lambda$ is the dimension of $\ker(\phi-\lambda \mathrm{id}_V)$.
\end{itemize}
Show that $\phi$ is diagonalizable if and only if $c_\phi(x)$ factors as a product of linear factors and for each eigenvalue $\lambda$ of $\phi$, the arithmetic multiplicity of $\lambda$ equals the geometric multiplicity of $\lambda$.
\end{problem}


\end{document}